\chapter{Kobordyzmy i rachunek uchwytów}
\label{ch:rachunek-uchwytow}

W Rozdziale~\ref{ch:teoria-morsea} punkt krytyczny
funkcji Morse'a dostarczał \emph{jednego} uchwytu.
Teraz interesuje nas cały układ uchwytów i jego
zmiany. Dwie różne listy uchwytów mogą opisywać tę
samą rozmaitość: można przestawiać ich kolejność,
przesuwać jeden po drugim, a czasem usunąć parę.
Właśnie takie ruchy stoją za twierdzeniem Smale'a
o $h$-kobordyzmie. Będziemy dokładnie odróżniać
\emph{geometryczne} zniknięcie pary od samego
skrócenia kompleksu łańcuchów.

W całym rozdziale rozmaitości są gładkie, zwarte
i bez naroży, chyba że chwilowo sklejamy je
wzdłuż kołnierza; naroża wtedy wygładzamy.
Zakładamy $m=\dim W\geq1$, aby wymiar brzegu wynosił
$n=m-1$. Obramowanie i otoczenie tubularne są już
w Sekcji~\ref{sec:obramowanie-normalne}, a pojedynczy
uchwyt $D^k\times D^{m-k}$ w
Definicji~\ref{def:uchwyt-morse-17}.

\section{Od funkcji Morse'a do rozkładu kobordyzmu}
\label{sec:kobordyzm-18}

\begin{definicja}[Kobordyzm]
\index{kobordyzm@Kobordyzm}
\label{def:kobordyzm-18}
\emph{Kobordyzmem} z $N_-^{m-1}$ do
$N_+^{m-1}$ nazywamy zwartą rozmaitość $W^m$
wraz z identyfikacją jej brzegu
$\partial W=N_-\sqcup N_+$.
Kołnierze brzegu pozwalają odczytać $N_-$
jako brzeg wejściowy, a $N_+$ jako wyjściowy.
Rozmaitości puste są dopuszczalne. Iloczyn
$N\times[0,1]$ jest przykładem kobordyzmu
z $N$ do $N$.
\end{definicja}

W \emph{kobordyzmie zorientowanym} identyfikacje brzegu
spełniają $\partial W=(-N_-)\sqcup N_+$,
gdzie orientację $\partial W$ wyznacza kolejność
„normalna zewnętrzna, baza styczna brzegu”.
Sama teoria uchwytów nie wymaga orientacji $W$.
Całkowite współczynniki brzegu kompleksu uchwytowego
również można zdefiniować bez niej: wystarczą orientacje
rdzeni i wynikające z nich koorientacje sfer pasa.

\begin{stwierdzenie}[Funkcja Morse'a zgodna z brzegiem]
\label{prop:morse-kobordyzm-18}
Istnieje funkcja Morse'a
$F:W\to[0,1]$ taka, że
$F^{-1}(0)=N_-$, $F^{-1}(1)=N_+$,
w kołnierzach nie ma punktów krytycznych,
a wszystkie punkty krytyczne leżą we wnętrzu
i mają różne wartości.
\end{stwierdzenie}
\begin{proof}
Twierdzenie o kołnierzu z
Sekcji~\ref{sec:kolnierz-15} daje rozłączne
obrazy $N_\pm\times[0,2\eta)$.
Na kołnierzu wejściowym kładziemy $F(x,t)=t$,
a na wyjściowym $F(x,t)=1-t$.
W szczególności różniczka $F$ w kierunku
normalnym jest tam niezerowa.
Funkcję przedłużamy gładko na środek $W$
z wartościami w $(\eta,1-\eta)$; można to
zrobić za pomocą podziału jedności, po
uprzednim ustaleniu wartości na nieco mniejszych
kołnierzach. Teraz stosujemy względną wersję
argumentu z Twierdzenia~\ref{thm:gestosc-morse}:
mała perturbacja wsparta \emph{wyłącznie}
w zwartym wnętrzu czyni przekrój $\dd F$
poprzecznym do przekroju zerowego $T^*W$.
W punkcie jego zera linearyzacja $\dd F$
jest hesjanem; poprzeczność oznacza jego
nieosobliwość. Zwartość daje skończenie wiele
takich punktów. Wybierając osobne małe
otoczenia punktów krytycznych i dodając
w każdym niewielką stałą z funkcją odcinającą,
rozdzielamy ich wartości. Perturbacja jest
stała blisko każdego punktu, a na dopełnieniu
tych otoczeń $|\dd F|$ ma dodatnią dolną
granicę; po dostatecznym zmniejszeniu zmiany
nowe punkty krytyczne nie powstają.
Kołnierze i wartości na końcach nie zmieniły się.
\end{proof}

\begin{twierdzenie}[Rozkład na uchwyty]
\label{thm:rozklad-uchwyty-18}
Niech $F$ będzie funkcją z
Stwierdzenia~\ref{prop:morse-kobordyzm-18}.
Jeżeli jej punkty krytyczne mają indeksy
$k_1,\ldots,k_s$ w kolejności rosnących
wartości, to
\begin{equation}\label{eq:rozklad-kobordyzm-18}
 W\cong
 (N_-\times[0,1])\cup H_{k_1}\cup\cdots\cup H_{k_s},
 \qquad H_k=D^k\times D^{m-k},
\end{equation}
z dokładnością do wygładzenia miejsc sklejenia.
Uchwyty są przyczepiane kolejno do aktualnego
brzegu wyjściowego. Dla $N_-=\varnothing$
zaczynamy od pustej przestrzeni.
\end{twierdzenie}
\begin{proof}
Wokół każdej wartości krytycznej $c_i$ wybierzmy
dwie wartości regularne $b_i^-<c_i<b_i^+$ tak, aby
przedziały $[b_i^-,b_i^+]$ były rozłączne i zawierały
po jednym punkcie krytycznym. Pozostałe pasy są regularne.
Małe otoczenie wejścia jest ustalonym kołnierzem;
gdy $N_-=\varnothing$, dolna podpoziomica jest pusta.
W każdym zwartym paśmie regularnym
$\nabla F/|\nabla F|^2$ jest gładkie i zmienia $F$
z szybkością $1$. Przepływ identyfikuje pas
z poziomicą razy przedział, jak w
Twierdzeniu~\ref{thm:pas-bez-krytycznych-15}.
Na paśmie zawierającym jeden punkt krytyczny
stosujemy Twierdzenie~\ref{thm:uchwyt-morse-17}:
jego dowód jest lokalny przy punkcie, a poza nim
korzysta tylko ze zwartości regularnego dopełnienia.
Otrzymujemy jeden $H_{k_i}$, z jego częścią przyczepianą
w aktualnym brzegu. Kolejne kołnierze wchłaniamy
przez zmianę parametru przedziału.
To daje~\eqref{eq:rozklad-kobordyzm-18}, także przy
pustej liście punktów krytycznych: wówczas cały
kobordyzm jest iloczynem. Nie zakładamy, że każda
składowa $W$ spotyka $N_-$; nowe składowe rozpoczynają
się od uchwytów indeksu $0$.
\end{proof}

\begin{przyklad}[Walec i dysk]
Na $N\times[0,1]$ funkcja $F(x,t)=t$
nie ma punktów krytycznych: rozkład nie
potrzebuje uchwytów. Dysk $D^m$, traktowany
jako kobordyzm z pustego zbioru do
$S^{m-1}$, ma na dysku jednostkowym funkcję
$F(x)=(1+|x|^2)/2$. Jej minimum to $1/2$,
więc $F^{-1}(0)=\varnothing$, a $F^{-1}(1)=S^{m-1}$.
W kołnierzu o współrzędnej $t=(1-|x|^2)/2$
mamy $F=1-t$. Jedyny punkt krytyczny w środku
ma indeks $0$.
To jeden $0$-uchwyt.
\end{przyklad}

\section{Anatomia uchwytu i rozkład dualny}
\label{sec:anatomia-uchwytu-18}

Zapiszmy $H_k=D^k_u\times D^{m-k}_v$.
Przyczepiamy go wzdłuż
$A=S^{k-1}\times D^{m-k}$,
więc samo wskazanie sfery $S^{k-1}$
nie wystarcza: trzeba także wskazać,
jak biegnie jej normalne pogrubienie.
To właśnie obramowanie z
Definicji~\ref{def:obramowanie-normalne}.

\begin{definicja}[Rdzeń, współrdzeń i dwie sfery]
\index{rdzen wspolrdzen i dwie sfery@Rdzeń, współrdzeń i dwie sfery}
\label{def:rdzen-pas-18}
\emph{Rdzeniem} jest $D^k\times\{0\}$,
a \emph{współrdzeniem} $\{0\}\times D^{m-k}$.
Brzeg rdzenia
$S^{k-1}\times\{0\}$ to \emph{sfera przyczepienia}.
Brzeg współrdzenia
$\{0\}\times S^{m-k-1}$ w \emph{nowym}
brzegu kobordyzmu to \emph{sfera pasa}
(ang.\ \emph{belt sphere}).
Dla $k=0$ sfera przyczepienia jest pusta;
dla $k=m$ sfera pasa jest pusta.
\end{definicja}

Sfera pasa jest miejscem, w którym kolejny
uchwyt może przeciąć poprzedni. Jeśli indeksy
dwóch kolejnych uchwytów wynoszą $k$ i $k+1$,
to ich sfery mają w pośrednim brzegu wymiaru
$m-1$ wymiary odpowiednio $m-k-1$ oraz $k$.
Suma tych wymiarów wynosi $m-1$, więc
przecięcie poprzeczne składa się z punktów.

\begin{figure}[htbp]
\centering
\begin{tikzpicture}[>=Latex,line cap=round,font=\small,scale=1.02]
 \shade[top color=manifoldblue!18,bottom color=manifoldblue!54]
      (-2.15,-.57) rectangle (2.15,.57);
 \draw[manifoldblue!75!black,thick]
      (-2.15,.57)--(2.15,.57)
      (-2.15,-.57)--(2.15,-.57);
 \fill[manifoldgold!39] (-2.15,0) ellipse (.25 and .57);
 \draw[manifoldgold!78!black,thick]
      (-2.15,0) ellipse (.25 and .57);
 \fill[manifoldgold!39] (2.15,0) ellipse (.25 and .57);
 \draw[manifoldgold!78!black,thick]
      (2.15,0) ellipse (.25 and .57);
 \draw[manifoldred!85!black,dashed,very thick,->]
      (-2.15,0)--(2.15,0);
 \fill[manifoldgreen!27,opacity=.82]
      (0,0) ellipse (.25 and .57);
 \draw[manifoldgreen!70!black,dashed,thick]
      (0,.57) arc[start angle=90,end angle=270,
                    x radius=.25,y radius=.57];
 \draw[manifoldgreen!70!black,very thick]
      (0,-.57) arc[start angle=270,end angle=450,
                    x radius=.25,y radius=.57];
 \node[anchor=south] at (-2.15,.79) {$\{-1\}\times D^2$};
 \node[anchor=south] at (2.15,.79) {$\{1\}\times D^2$};
 \node[manifoldgreen!65!black,anchor=south]
      at (0,1.27) {współrdzeń $D^2$};
 \draw[->,manifoldgreen!65!black]
      (0,1.20)--(.16,.42);
 \node[manifoldred!75!black] at (0,-1.15)
      {rdzeń $D^1$};
 \node[manifoldgreen!65!black,anchor=west]
      at (2.62,-.35) {sfera pasa $S^1$};
 \draw[->,manifoldgreen!65!black]
      (2.60,-.47)--(.177,-.403);
\end{tikzpicture}
\caption{Uchwyt indeksu $1$ w trójwymiarowym
kobordyzmie: dwie złote podstawy są przyczepiane
do starego brzegu. Czerwona przerywana oś jest
rdzeniem; zielony środkowy dysk jest współrdzeniem,
a jego okrąg brzegowy to sfera pasa. Rysunek
pokazuje sam uchwyt, bez reszty kobordyzmu.}
\label{fig:anatomia-uchwytu-18}
\end{figure}

\begin{twierdzenie}[Uchwyt widziany od drugiej strony]
\label{thm:dualny-uchwyt-18}
Po odwróceniu kierunku kobordyzmu uchwyt
indeksu $k$ w wymiarze $m$ staje się uchwytem
indeksu $m-k$. Jego nowym rdzeniem jest stary
współrdzeń, a nową sferą przyczepienia stara
sfera pasa. Cała lista indeksów odwraca kolejność
i zamienia $k$ na $m-k$.
\end{twierdzenie}
\begin{proof}
W modelu $F=c-|u|^2+|v|^2$ z
Lematu~\ref{thm:lemat-morsea-17}
odwrócenie funkcji daje
$-F=-c-|v|^2+|u|^2$.
Ujemna podprzestrzeń hesjanu dla $-F$
to dotychczasowa dodatnia część wymiaru
$m-k$. Zamiana miejsc czynników w
$D^k_u\times D^{m-k}_v$ daje
$D^{m-k}_v\times D^k_u$.
Jej część przyczepiana leży w brzegu,
który wcześniej był wyjściowy:
jest to $S^{m-k-1}\times D^k$.
Stąd nowa sfera przyczepienia to
$S^{m-k-1}\times\{0\}$, czyli stara
sfera pasa. Przechodzenie poziomic
od $N_+$ do $N_-$ odwraca także
porządek wszystkich punktów krytycznych.
\end{proof}

\section{Co dzieje się z brzegiem: ślad chirurgii}
\label{sec:slad-chirurgii-18}

W rozdziale poprzednim patrzyliśmy na całe
podpoziomice. Teraz chcemy zrozumieć tylko
ich \emph{brzegi}, czyli rozmaitości $N_t=F^{-1}(t)$.
Zmiana brzegu jest operacją chirurgiczną
z Definicji~\ref{def:chirurgia-zarys-16}.

\begin{twierdzenie}[Uchwyt jako ślad chirurgii]
\label{thm:slad-chirurgii-18}
Niech $1\leq k\leq m-1$ i załóżmy, że
elementarny kobordyzm $W$ powstaje z
$N_-^{m-1}\times[0,1]$ przez dołączenie
jednego $k$-uchwytu za pomocą
\[
 \alpha:S^{k-1}\times D^{m-k}
       \hookrightarrow N_-\times\{1\}.
\]
Po wygładzeniu naroży jego drugi brzeg wynosi
\begin{equation}\label{eq:brzeg-chirurgia-18}
 N_+\cong
 \bigl(N_-\setminus\operatorname{int}
        \alpha(S^{k-1}\times D^{m-k})\bigr)
 \mathop{\cup}_{S^{k-1}\times S^{m-k-1}}
       (D^k\times S^{m-k-1}).
\end{equation}
Jest to chirurgia indeksu $k-1$ na
$(m-1)$-wymiarowym $N_-$.
Odwrotnie, każda taka obramowana chirurgia
ma ślad będący elementarnym kobordyzmem
z jednym $k$-uchwytem.
\end{twierdzenie}
\begin{proof}
Brzeg uchwytu ma dwa kawałki:
\[
 \partial(D^k\times D^{m-k})
 =\underbrace{S^{k-1}\times D^{m-k}}_{A}
   \cup
   \underbrace{D^k\times S^{m-k-1}}_{B}.
\]
Kawałki spotykają się dokładnie we wspólnym
brzegu $S^{k-1}\times S^{m-k-1}$.
Podczas przyklejania $A$ zostaje utożsamiony
z $\alpha(A)$ leżącym w starym brzegu;
punkty wnętrza $A$ przestają należeć do brzegu całości.
Punkty $\partial A=\partial B$ leżą jeszcze na narożu,
które następnie wygładzamy.
Pozostała część $N_-$ nadal jest w brzegu,
a kawałek $B$ staje się \emph{nową} częścią
brzegu. To daje~\eqref{eq:brzeg-chirurgia-18}.
Żaden inny fragment $\partial H_k$ nie pozostaje.
Identyfikacja $\alpha$ obejmuje normalne dyski
$D^{m-k}$, więc zawiera dane obramowania;
bez nich wyrażenie po prawej nie jest
jednoznaczną operacją gładką.

W drugą stronę bierzemy $N_-\times[0,1]$
i przyklejamy $D^k\times D^{m-k}$ do górnego
brzegu przez zadane obramowane zanurzenie
$\alpha$. To buduje zwarty kobordyzm, którego
brzeg wyjściowy jest dokładnie prawą stroną
\eqref{eq:brzeg-chirurgia-18}. Wygładzenie
naroży odbywa się w kołnierzu wspólnego brzegu.
\end{proof}

\begin{uwaga}[Skrajne indeksy]
Przy $k=0$ nic nie przyklejamy do starego
brzegu: powstaje nowa składowa $D^m$,
której brzeg to $S^{m-1}$. Przy $k=m$
uchwyt przykleja się całym $S^{m-1}$,
zamykając odpowiednią składową brzegu.
Dlatego w~\eqref{eq:brzeg-chirurgia-18}
ograniczyliśmy się do $1\leq k\leq m-1$.
\end{uwaga}

\begin{przyklad}[Sfera przechodzi w torus]
\label{prz:sfera-torus-chirurgia-18}
Weźmy $m=3$, $k=1$ i $N_-=S^2$.
Wybierzmy przyklejenie zgodne z orientacją,
tak aby orientacja $S^2\times[0,1]$ przedłużała się
na dołączony uchwyt.
Obszar $S^0\times D^2$ to dwa rozłączne
dyski w sferze. Po usunięciu ich wnętrz
pozostaje pierścień $S^1\times[0,1]$.
Nowa część brzegu $D^1\times S^1$
jest \emph{drugim} pierścieniem; sklejenie
obu wzdłuż dwóch okręgów daje spójną, zorientowaną
powierzchnię o charakterystyce Eulera $0$, czyli $T^2$.
Identyfikację z torusem można uzasadnić bez klasyfikacji powierzchni.
Po sklejeniu jednej pary okręgów mamy dłuższy pierścień;
sklejenie jego końców opisuje dyfeomorfizm $h:S^1\to S^1$.
Orientowalność oznacza, że $h$ zachowuje orientację okręgu.
Jego podniesienie $\widetilde h:\R\to\R$ spełnia
$\widetilde h(x+1)=\widetilde h(x)+1$ oraz $\widetilde h'(x)>0$.
Dlatego $(1-t)x+t\widetilde h(x)$, $0\leq t\leq1$,
opuszcza się do izotopii $h_t$ od identyczności do $h$.
Wybierając parametr tej izotopii stały blisko końców,
otrzymujemy dyfeomorfizm
\[
 (S^1\times[0,1])/((x,1)\sim(h(x),0))\longrightarrow S^1\times S^1,
 \qquad [x,t]\longmapsto(h_t(x),e^{2\pi it}).
\]
Stałość izotopii przy końcach zapewnia gładkość po sklejeniu.
Warunek orientacji jest istotny: przy drugim typie
sklejenia pierścieni powstaje butelka Kleina,
a trójwymiarowy kobordyzm jest nieorientowalny.
Jeśli dodatkowo wypełnimy wejściowy brzeg
kulą $D^3$, powstanie $D^3$ z dołączonym
$1$-uchwytem, czyli pełny torus $D^2\times S^1$.
Sam kobordyzm ma \emph{dwa} brzegi: $S^2$ i $T^2$.
Zgodność charakterystyk Eulera kontroluje
wynik: $\chi(S^2)=2$, usuwamy dwa dyski
o łącznej charakterystyce $2$ i wklejamy
pierścień o charakterystyce $0$, więc
$\chi(T^2)=0$. To sprawdzenie nie zastępuje
opisu sposobu sklejenia.
\end{przyklad}

\begin{figure}[htbp]
\centering
\begin{tikzpicture}[>=Latex,line cap=round,font=\small]
 % Lewy brzeg: sfera z dwoma zaznaczonymi dyskami.
 \shade[inner color=white,outer color=manifoldblue!55]
      (-3.12,0) circle (1.03);
 \draw[manifoldblue!68!black,thick] (-3.12,0) circle (1.03);
 \shade[inner color=manifoldgold!15,outer color=manifoldgold!63]
      (-3.54,-.04) ellipse (.21 and .33);
 \shade[inner color=manifoldgold!15,outer color=manifoldgold!63]
      (-2.70,-.04) ellipse (.21 and .33);
 \draw[manifoldgold!78!black,thick]
      (-3.54,-.04) ellipse (.21 and .33)
      (-2.70,-.04) ellipse (.21 and .33);
 \node at (-3.12,-1.70) {$S^2$: dwa dyski};
 \draw[->,very thick,manifoldred!75!black]
       (-1.81,0)--(-.60,0);
 \node[align=center,manifoldred!75!black] at (-1.19,.75)
       {$1$-uchwyt};
 % Prawy brzeg: gładki model parametryczny torusa.
 \begin{scope}[shift={(1.60,0)}]
   \begin{axis}[hide axis,view={0}{34},axis equal image,scale=2,
      xmin=-3,xmax=3,ymin=-3,ymax=3,zmin=-.8,zmax=.8,scale only axis,
      at={(0,0)},anchor=center,width=9cm,height=5cm,
       colormap={toruslight}{color(0cm)=(manifoldblue!67);
         color(.42cm)=(manifoldblue!33);color(1cm)=(white)}]
    \addplot3[surf,shader=interp,draw=none,z buffer=sort,
      samples=34,samples y=68,domain=0:360,y domain=0:360,
      variable=\t,variable y=\p]
     ({(2+0.69*cos(\t))*cos(\p)},
      {(2+0.69*cos(\t))*sin(\p)},
      {0.69*sin(\t)});
   \end{axis}
 \end{scope}
 \node at (1.60,-1.70) {$T^2$: nowy brzeg};
\end{tikzpicture}
\caption{Dołączenie $1$-uchwytu zgodnie z orientacją zmienia
brzeg wyjściowy trójwymiarowego kobordyzmu z $S^2$ na $T^2$.
Zaznaczone na sferze złote dyski stanowią
$S^0\times D^2$; po wycięciu ich wnętrz
doklejamy pierścień $D^1\times S^1$.}
\label{fig:chirurgia-sfera-torus-18}
\end{figure}

\begin{przyklad}[Operacja odwrotna i dualny uchwyt]
Odwróćmy poprzedni kobordyzm. Ten sam
trójwymiarowy kawałek widzimy teraz jako
$2$-uchwyt przyczepiony do torusa.
W pośrednim $T^2$ sfera pasa wcześniejszego
$1$-uchwytu jest okręgiem \emph{południkowym}.
Przy odwróceniu kobordyzmu właśnie ona staje się
sferą przyczepienia dualnego $2$-uchwytu,
zgodnie z Twierdzeniem~\ref{thm:dualny-uchwyt-18}.
Wzór~\eqref{eq:brzeg-chirurgia-18}
opisuje przejście $T^2\to S^2$ jako
chirurgię indeksu $1$. Obramowanie tej
krzywej jest tutaj odziedziczone po pierwszym uchwycie.
Ogólnie przyklejenie określa obramowane pogrubienie,
a nie tylko obraz sfery przyczepienia.
\end{przyklad}

\section{Łańcuchy uchwytów i liczby przecięcia}
\label{sec:lancuchy-uchwytow-18}

Uchwyt jest pogrubioną komórką. Z tego powodu
rozkład~\eqref{eq:rozklad-kobordyzm-18} powinien
prowadzić do kompleksu obliczającego homologię
\emph{względną} $H_*(W,N_-)$. Ważna nowość
polega na tym, że jego różniczkę można odczytać
z przecięcia dwóch sfer w poziomicy.

\begin{twierdzenie}[Kompleks względny uchwytów]
\label{thm:kompleks-uchwytow-18}
Załóżmy, że uchwyty ułożono według indeksów,
tak że wszystkie uchwyty indeksu $k-1$
przyczepiono przed uchwytami indeksu $k$.
Taki porządek można osiągnąć ruchem opisanym
w Twierdzeniu~\ref{thm:przestawianie-18}.
Po wybraniu orientacji rdzeni istnieje kompleks
\[
 C_k(W,N_-)=\Z\langle
       \text{uchwyty indeksu }k\rangle,
 \qquad H_k(C_\bullet,\partial)\cong H_k(W,N_-;\Z).
\]
Jeśli $p$ jest $k$-uchwytem, a $q$ jest
$(k-1)$-uchwytem, to współczynnik przy $q$
w $\partial_k p$ jest podpisaną liczbą
poprzecznych przecięć w pośredniej poziomicy:
\begin{equation}\label{eq:macierz-uchwyty-18}
 \langle\partial_kp,q\rangle
 = I\bigl(A_p,B_q\bigr),\qquad
 A_p=S^{k-1}\times\{0\},\quad
 B_q=\{0\}\times S^{m-k}.
\end{equation}
Sfery przenosimy do tej samej poziomicy
przez regularne pasy przepływem.
\end{twierdzenie}
\begin{proof}
Niech $W^{(k)}$ będzie częścią po dołączeniu
uchwytów indeksów co najwyżej $k$,
a $W^{(-1)}=N_-$ wraz z małym kołnierzem.
Wycięcie pozostałej przestrzeni i
ściągnięcie drugiego czynnika każdego
$D^k\times D^{m-k}$ dają
\[
 H_j(W^{(k)},W^{(k-1)};\Z)
 \cong
 \bigoplus_{p:\,\lambda(p)=k}
 H_j(D^k,S^{k-1};\Z)
 =
 \begin{cases}
 \Z^{c_k},&j=k,\\
 0,&j\ne k.
 \end{cases}
\]
Dla uchwytów tego samego indeksu wybieramy pole
Morse'a--Smale'a; nie ma między nimi trajektorii.
Ich sfery wyjściowe można więc sprowadzić do jednego
poziomu poniżej całego bloku. Czasy trafienia są
gładkie i ograniczone na tych zwartych sferach;
inaczej granica dałaby trajektorię do punktu
o tym samym indeksie. Sfery są rozłączne, gdyż
trajektoria ma tylko jeden początek krytyczny.
Ich dostatecznie małe pogrubienia też są rozłączne.
Pozwala to dołączyć cały blok jednocześnie
i uzasadnia powyższą sumę prostą przez wycięcie
po uprzednim pogrubieniu części przyczepianych kołnierzami.
Orientacja dysku rdzenia wybiera generator
odpowiedniego składnika $\Z$.

Różniczka jest złożeniem
\[
 H_k(W^{(k)},W^{(k-1)})
 \xrightarrow{\delta}H_{k-1}(W^{(k-1)},N_-)
 \longrightarrow H_{k-1}(W^{(k-1)},W^{(k-2)}).
\]
Pierwsze odwzorowanie pochodzi z ciągu dokładnego trójki
$(W^{(k)},W^{(k-1)},N_-)$, a drugie zapomina
część względną. W stopniu $0$ brzeg jest zerowy.
Jego kwadrat jest zerowy: brzeg brzegu
rdzenia znika, co można też sprawdzić
w długich ciągach dokładnych kolejnych
potrójek. Ponieważ każda para
$(W^{(k)},W^{(k-1)})$ ma homologię
skupioną w jednym stopniu, kolejne
długie ciągi dokładne, albo równoważny
kompleks CW z
Twierdzenia~\ref{thm:cw-morse-17},
identyfikują homologię tego kompleksu
z $H_*(W,N_-)$.

Pozostaje uzasadnić konkretny wzór.
Brzeg zorientowanego rdzenia $p$
to jego sfera przyczepienia $A_p$.
Przenosimy ją do regularnej poziomicy,
w której widoczny jest współrdzeń $q$.
Po zgnieceniu innych komórek niż $q$
współczynnik brzegowy jest stopniem
odwzorowania
$S^{k-1}\to S^{k-1}$.
Wybieramy wartość regularną w komórce $q$.
Jej przeciwobrazy odpowiadają dokładnie
przecięciom $A_p$ ze sferą pasa $B_q$:
współrdzeń $q$ jest poprzecznym dyskiem,
a jego brzeg leży w rozważanej poziomicy.
Orientujemy $A_p$ jako brzeg rdzenia $p$.
Normalną orientację $B_q$ wyznacza zorientowany
rdzeń $q$: w modelu uchwytu jego kierunki są
normalne do współrdzenia. Znak przecięcia to znak
izomorfizmu $T_zA_p\to\nu_z B_q$ w poziomicy.
Jest on znakiem lokalnego stopnia. Ta konwencja
działa także dla nieorientowalnego $W$ i zgadza się
z konwencją w dowodzie
Twierdzenia~\ref{thm:homologia-morse-17}.
Sumując znaki, otrzymujemy
\eqref{eq:macierz-uchwyty-18}.
Zmiana orientacji rdzenia zmienia tylko
znak odpowiedniego generatora.
\end{proof}

\begin{przyklad}[Jeden punkt przecięcia lub żaden]
\label{prz:meridian-longitude-18}
Zacznijmy od $S^2\times[0,1]$ i
dołączmy $1$-uchwyt, otrzymując na górze $T^2$.
Jego sfera pasa $B_1$ jest południkiem
torusa. Jeżeli kolejny $2$-uchwyt przyczepimy
wzdłuż krzywej biegnącej raz w kierunku
równoleżnikowym, to sfera przyczepienia
$A_2$ przecina $B_1$ w jednym punkcie.
Po dobraniu orientacji
$C_2=\Z\xrightarrow{\ \partial_2=\pm1\ }
 C_1=\Z$.
Kompleks względny jest acykliczny,
a geometryczne zniesienie pary
udowodnimy w
Twierdzeniu~\ref{thm:zniesienie-18}.

Jeśli natomiast przyczepimy $2$-uchwyt
wzdłuż krzywej \emph{równoległej do
południka} $B_1$, po małej perturbacji
krzywe te są rozłączne i $\partial_2=0$.
To inny wybór niż w parze znoszącej się.
W odwróconym pojedynczym kobordyzmie
z poprzedniego przykładu południk
jest właśnie sferą przyczepienia
\emph{dualnego} $2$-uchwytu;
nie należy mylić tych dwóch sytuacji.
\end{przyklad}

\begin{figure}[htbp]
\centering
\begin{tikzpicture}[>=Latex,font=\small,line cap=round,scale=1.15]
 \fill[manifoldblue!6] (-1.45,-1.15) rectangle (1.45,1.15);
 \draw[manifoldblue!70!black,thick]
      (-1.45,-1.15) rectangle (1.45,1.15);
 \draw[manifoldgreen!70!black,very thick,->]
      (0,-1.15)--(0,.85);
 \draw[manifoldgreen!70!black,very thick]
      (0,.85)--(0,1.15);
 \draw[manifoldred!82!black,very thick,->]
      (-1.45,0)--(1.18,0);
 \draw[manifoldred!82!black,very thick]
      (1.18,0)--(1.45,0);
 \fill[manifoldgold!85!black] (0,0) circle (2.5pt);
 \node[manifoldgreen!65!black,anchor=west] at (.12,.78)
       {$B_1$};
 \node[manifoldred!80!black,anchor=south] at (.95,.09)
       {$A_2$};
 \node[anchor=north,align=center] at (0,-1.35)
       {przeciwległe boki są sklejone: $T^2$};
\end{tikzpicture}
\caption{Pośredni torus przedstawiony jako kwadrat
ze sklejonymi bokami. Zielona sfera pasa
$B_1$ i czerwona sfera przyczepienia
$A_2$ przecinają się raz. Dlatego odpowiedni
współczynnik macierzy $\partial_2$ wynosi
$+1$ albo $-1$, zależnie od orientacji.}
\label{fig:przeciecie-uchwytow-18}
\end{figure}

\begin{uwaga}[Rachunek algebraiczny to jeszcze nie geometria]
Jeżeli współczynnik
$\langle\partial_{k+1}p,q\rangle=\pm1$,
to po zmianach baz można wydzielić podkompleks
$0\to\Z\xrightarrow{\pm1}\Z\to0$ i go usunąć.
Najpierw zerujemy pozostałe wpisy w wierszu i kolumnie
tego współczynnika. Tożsamość $\partial^2=0$
zeruje wówczas odpowiednie wpisy w sąsiednich
różniczkach, więc naprawdę otrzymujemy sumę prostą
kompleksów. Usunięty składnik jest ściągalny.
Nie wynika z tego automatycznie, że
sfery $A_p$ i $B_q$ mają \emph{jeden
geometryczny} punkt przecięcia:
podpisana suma może wynosić $1$
przy trzech, pięciu albo większej
liczbie punktów. Ich usunięcie wymaga
argumentu geometrycznego, często
triku Whitneya i warunków wymiarowych.
\end{uwaga}

\section{Przestawianie i przesuwanie uchwytów}
\label{sec:ruchy-uchwytow-18}

Lista uchwytów zależy od funkcji Morse'a,
choć sam kobordyzm pozostaje ten sam.
Dwa podstawowe ruchy mają różne znaczenie:
\emph{przestawienie} zmienia kolejność
wartości krytycznych, a \emph{przesunięcie}
(ang.\ \emph{handle slide}) zmienia sposób
przyczepienia uchwytów tego samego indeksu.

\begin{definicja}[Pole gradientopodobne spadku]
\index{pole gradientopodobne spadku@Pole gradientopodobne spadku}
\label{def:pole-spadku-18}
Przez pole gradientopodobne dla $F$ rozumiemy pole $X$
spełniające $XF<0$ poza punktami krytycznymi,
zerujące się w tych punktach i będące blisko każdego
z nich ujemnym gradientem $F$ dla pewnej gładkiej metryki.
Można wybrać $X=-\nabla_g F$ na całym kobordyzmie,
z metryką euklidesową w małych mapach Morse'a.
W kołnierzach $X$ jest skierowane od wyjścia do wejścia.
\end{definicja}

Poza zerami takie pole również można zrealizować
jako ujemny gradient: ustalamy
$g(X,X)=-\dd F(X)>0$, $g(X,\ker\dd F)=0$
i dodatnią metrykę na $\ker\dd F$.
Podział jedności skleja te metryki z ustalonymi
przy zerach, zachowując $g(X,\cdot)=-\dd F$.
W konsekwencji wcześniejsze argumenty o końcach
trajektorii stosują się aż do chwili trafienia w brzeg.

\begin{twierdzenie}[Przestawianie dwóch punktów]
\label{thm:przestawianie-18}
Niech $p,q$ będą jedynymi punktami krytycznymi
pewnego pasa, $F(p)<F(q)$.
Jeżeli dla wybranego gradientopodobnego pola
spadku nie ma trajektorii $q\to p$,
można zmienić $F$ we wnętrzu tego pasa
bez tworzenia nowych punktów krytycznych,
tak aby $p$ i $q$ zamieniły kolejność wartości.
Pole $X$ można zachować na \emph{całym} paśmie:
$X\widetilde F<0$ poza jego zerami, a blisko $p,q$
nowa funkcja różni się od starej tylko o stałe.
Przy brzegach pasa funkcja pozostaje niezmieniona.
W szczególności
rozkład uchwytowy można uporządkować
według niemalejących indeksów.
\end{twierdzenie}
\begin{proof}
Niech pas ma regularne końce $a<b$ i wybierzmy
$F(p)<s<F(q)$, $L=F^{-1}(s)$.
Zbiory $A=W^s(p)\cap L$ i $B=W^u(q)\cap L$
są rozłączne: punkt wspólny dawałby trajektorię $q\to p$.
Są też zwarte. Na przykład mała sfera wejściowa
lokalnego dysku stabilnego $p$ po przepływie wstecz
dociera do $L$, ponieważ między jej poziomem a $s$
nie ma punktów krytycznych. Czas trafienia jest
gładki na zwartej sferze, a jej obraz to $A$.
Tak samo otrzymujemy $B$ z lokalnej sfery niestabilnej $q$.
Zbiór może być pusty przy skrajnym indeksie.

Wybierzmy gładką $\rho_L:L\to[0,1]$, równą $1$
w otoczeniu $A$ i $0$ w otoczeniu $B$.
Przedłużmy ją na całe pasmo, nadając tę samą wartość
wszystkim punktom jednej trajektorii przecinającej $L$.
Trajektorie poniżej $L$, które nie przecinają $L$
wstecz, muszą mieć początek w $p$; na nich kładziemy $\rho=1$.
Trajektorie powyżej $L$, które nie dochodzą do $L$
naprzód, muszą kończyć się w $q$; na nich kładziemy $\rho=0$.
Tak samo definiujemy wartości w samych $p,q$.

Sprawdźmy gładkość tego przedłużenia.
Tam, gdzie trajektoria przecina $L$ w skończonym czasie,
wynika ona z gładkości czasu trafienia.
W dostatecznie małym otoczeniu $p$ wszystkie trajektorie,
które dochodzą do $L$ wstecz, trafiają w obszar $\rho_L=1$.
Gdyby nie, ciąg punktów dążących do $p$ dawałby
po wybraniu podciągu punkt graniczny na $L$ poza tym obszarem.
Jego trajektoria musiałaby zbiegać do $p$:
w przeciwnym razie dochodziłaby do dolnego brzegu
w skończonym czasie, podobnie jak pobliskie trajektorie,
zanim mogłyby zbliżyć się do $p$.
To przeczyłoby temu, że cały $A$ leży w obszarze $\rho_L=1$.
Zatem $\rho$ jest identycznie $1$ blisko $p$.
Jest też stała w otoczeniu każdego dalszego punktu
$W^u(p)$, po przeniesieniu tego otoczenia przepływem.
Analogicznie $\rho=0$ blisko $q$ i wzdłuż $W^s(q)$.
Otrzymujemy gładką funkcję na całym paśmie z $X\rho=0$.

Wybierzmy rosnące dyfeomorfizmy $h_p,h_q:[a,b]\to[a,b]$,
równe identyczności przy końcach, takie że
\[
 h_p(F(p))>h_q(F(q)).
\]
W małym otoczeniu odpowiedniej wartości krytycznej
każdy $h$ może być translacją. Takie odwzorowania konstruujemy,
zadając dodatnią pochodną o odpowiednich całkach
na przedziałach po obu stronach tej wartości.
Nie próbujemy przestawić dwóch liczb jednym odwzorowaniem rosnącym:
używamy dwóch różnych odwzorowań na różnych rodzinach trajektorii.
Kładziemy
\[
 \widetilde F=\rho\,h_p(F)+(1-\rho)h_q(F).
\]
Wtedy
\[
 X\widetilde F=
 \bigl(\rho h_p'(F)+(1-\rho)h_q'(F)\bigr)XF<0
\]
poza $p,q$, więc nowe punkty krytyczne nie powstają.
Blisko $p$ lub $q$ dodaliśmy tylko stałą,
zachowując hesjany i indeksy. Przy końcach pasa
oba odwzorowania $h$ są identycznością, więc także $\widetilde F=F$.
To dowodzi przestawienia przy niezmienionym polu.

Dla ostatniego zdania wybieramy
pole Morse'a--Smale'a jak w
Twierdzeniu~\ref{thm:gen-ms-17}.
Wtedy trajektoria od $q$ do $p$
wymaga $\lambda(q)>\lambda(p)$
(Stwierdzenie~\ref{prop:wymiar-ms-17}).
Jeśli dwa sąsiednie punkty, wymienione od niższej
do wyższej wartości funkcji, mają indeksy
w kolejności malejącej lub równej,
takiej trajektorii nie ma i wolno je
zamienić. Kolejne zamiany, jak przy
porządkowaniu skończonej listy, ustawiają
indeksy niemalejąco.
\end{proof}

\begin{uwaga}
Warunek o trajektorii jest istotny.
Jeśli $q\to p$ istnieje, przepływ
łączy dwa lokalne modele i niezależne
przesunięcie ich wartości może
przestać być zgodne z kierunkiem
spadku. Powyższe twierdzenie nie
pozwala dowolnie przestawiać wszystkich
uchwytów.
\end{uwaga}

\begin{stwierdzenie}[Przesunięcie uchwytu]
\label{prop:slide-18}
Załóżmy, że $p$ i $q$ są dwoma uchwytami
tego samego indeksu $1\leq k\leq m-1$,
przyczepionymi w jednym bloku rozkładu.
Dla $k\geq2$ wybieramy obramowaną wstęgę od sfery
przyczepienia $p$ do małej równoległej kopii sfery $q$,
wyznaczonej przez obramowanie uchwytu $q$.
Wnętrze wstęgi omija wszystkie sfery i ich używane
pogrubienia. Sferę $p$ zastępujemy sumą połączoną
z tą kopią; sam uchwyt $q$ pozostaje na miejscu.
Zakładamy istnienie takiej wstęgi z uzgodnionym
obramowaniem, nie dowolnego łuku przecinającego inne dane.
Dla $k=1$ przesuwamy jedną ze stóp
uchwytu $p$ przez uchwyt $q$.
Cały kobordyzm pozostaje dyfeomorficzny
względem brzegu wejściowego $N_-$.
W kompleksie z
Twierdzenia~\ref{thm:kompleks-uchwytow-18}
ruch ten jest elementarną zmianą bazy
\begin{equation}\label{eq:slide-baza-18}
 e_p'=e_p\pm e_q,\qquad e_q'=e_q.
\end{equation}
\end{stwierdzenie}
\begin{proof}
Przyczepmy najpierw uchwyt $q$.
W jego nowym brzegu $D^k\times S^{m-k-1}$
wybierzmy $v_0\in S^{m-k-1}$.
Równoległa kopia $S^{k-1}\times\{v_0\}$ dawnej
sfery przyczepienia ogranicza teraz dysk
$D^k\times\{v_0\}$ leżący \emph{w brzegu}.
Dysk ten jest równoległy do rdzenia; nie prowadzimy
izotopii po rdzeniu leżącym we wnętrzu kobordyzmu.
Dla $k\geq2$ sumę połączoną z brzegiem tego dysku
można zsunąć przez dysk na pierwotną sferę $p$.
W jego produktowym otoczeniu jest to zwykłe
przesunięcie paska, wraz z normalnym pogrubieniem.
Założenie o obramowanej wstędze zapewnia zgodność
tego pogrubienia po obu stronach.
Dla $k=1$ dysk jest odcinkiem i przesuwamy
jedną stopę uchwytu wzdłuż tego odcinka.

Jest to izotopia obramowanego przyklejenia $p$
w brzegu przestrzeni z już dołączonym $q$.
Twierdzenie o przedłużaniu izotopii z
Sekcji~\ref{sec:izotopie-otoczen-15} rozszerza ją
na otoczenie brzegu, z nośnikiem poza wcześniejszym
brzegiem wejściowym. Przedłużenie na uchwyt $p$
daje dyfeomorfizm dwóch rozkładów względem $N_-$.
Nie musi on punktowo unieruchamiać całej wcześniej
dołączonej części wraz z jej brzegiem wyjściowym.
Jeśli dołączono dalsze uchwyty, przenosimy także
ich odwzorowania przyczepienia przez ten dyfeomorfizm.

Prześledźmy rdzeń $D^k$ uchwytu $p$
przez ten ruch. Jego nowa klasa
w grupie względnej
$H_k(W^{(k)},W^{(k-1)})$
biegnie raz przez dawny rdzeń $p$
i raz, ze znakiem zależnym od
orientacji wstęgi, przez rdzeń $q$.
Pozostałe rdzenie nie są dotknięte.
To daje~\eqref{eq:slide-baza-18}.
Ponieważ $\partial$ jest liniowy,
$\partial_k e_p'
 =\partial_k e_p\pm\partial_k e_q$.
Jeśli występują uchwyty indeksu $k+1$,
ich współrzędne w stopniu $k$ zmieniają
się przez odwrotną zmianę bazy;
homologia pozostaje ta sama.
\end{proof}

\begin{figure}[htbp]
\centering
\begin{tikzpicture}[>=Latex,font=\small,line cap=round,scale=1.03]
 \draw[manifoldblue!85!black,very thick]
      (-3.35,0) circle (.56);
 \draw[manifoldgold!85!black,very thick]
      (-1.67,0) circle (.56);
 \draw[manifoldgold!75!black,densely dashed]
      (-1.67,0) circle (.68);
 \fill[manifoldblue!15]
      (-2.80,-.12)--(-2.34,-.12)
      --(-2.34,.12)--(-2.80,.12)--cycle;
 \draw[manifoldblue!65!black,densely dashed]
      (-2.80,-.12)--(-2.34,-.12)
      (-2.80,.12)--(-2.34,.12);
 \node at (-3.35,-.90) {$A_p$};
 \node at (-1.67,-.90) {$A_q$};
 \draw[->,gray!70,thick] (-.60,0)--(.45,0);
 \draw[manifoldred!85!black,very thick]
      (.85,0) .. controls (.85,.75) and (1.45,.90)
      .. (1.95,.40) .. controls (2.10,.18) and (2.10,.18)
      .. (2.25,.40) .. controls (2.70,.90) and (3.45,.75)
      .. (3.45,0) .. controls (3.45,-.75) and (2.70,-.90)
      .. (2.25,-.40) .. controls (2.10,-.18) and (2.10,-.18)
      .. (1.95,-.40) .. controls (1.45,-.90) and (.85,-.75)
      .. (.85,0);
 \draw[manifoldgold!85!black,very thick]
      (2.78,0) circle (.40);
 \node at (1.83,-.98) {$A_p\# A_q'$};
 \node at (3.15,-.98) {$A_q$};
\end{tikzpicture}
\caption{Schemat przesunięcia w przypadku,
gdy sfery przyczepienia są okręgami.
Po lewej wybieramy wstęgę między
$A_p$ i przerywaną równoległą kopią $A_q'$ sfery $A_q$;
sam złoty okrąg $A_q$ pozostaje na miejscu.
Po prawej nowa sfera
przyczepienia pierwszego uchwytu
przebiega raz także wzdłuż drugiego.
Dokładne osadzenie i obramowanie
w wyższym wymiarze nie są tu pokazane.}
\label{fig:slide-18}
\end{figure}

\section{Kiedy dwa uchwyty naprawdę się znoszą?}
\label{sec:znoszenie-uchwytow-18}

Na poziomie algebry różniczka $1:\Z\to\Z$
ma jądro i koker zerowe. Geometrycznie
potrzebujemy czegoś mocniejszego: dwie
odpowiednie sfery muszą spotkać się
\emph{w jednym punkcie}, a nie tylko
z algebraiczną sumą znaków równą $1$.

\begin{twierdzenie}[Geometryczne zniesienie pary]
\label{thm:zniesienie-18}
Niech pas kobordyzmu zawiera dokładnie
dwa punkty krytyczne $p,q$ o indeksach
odpowiednio $k,k+1$ i wartościach
$F(p)<F(q)$. Załóżmy, że dla wybranego
pola gradientopodobnego sfera przyczepienia
$A_q$ przecina sferę pasa $B_p$ w jednej
regularnej poziomicy \emph{poprzecznie
w dokładnie jednym punkcie}.
Wtedy można zmienić funkcję Morse'a
we wnętrzu pasa, nie ruszając jej przy
jego brzegach, tak by nie miała w nim
punktów krytycznych. Pas jest zatem
iloczynowym kobordyzmem: para uchwytów
$H_k,H_{k+1}$ znosi się.
\end{twierdzenie}
\begin{proof}
\emph{Krok 1: jedna trajektoria.}
Każda trajektoria z $q$ do $p$ przecina
każdą regularną poziomicę między ich
wartościami dokładnie raz, bo $F$
ściśle maleje poza punktami krytycznymi.
Punkt przecięcia $A_q\cap B_p$ wyznacza
taką trajektorię. Przy jednym punkcie
przecięcia jest więc dokładnie jedna
trajektoria łącząca, oznaczmy ją $\gamma$.
Jej przecięcie jest poprzeczne, zatem
nie znika przy małej zmianie przekroju.

\emph{Krok 2: dostosowujemy modele przy końcach.}
Nie zakładamy, że dowolny gradient da się gładko
zlinearyzować. Wolno nam natomiast zmieniać pole
spadku, zachowując funkcję i potrzebne dyski.
Przy $q$ zachowamy lokalny dysk $W^u(q)$,
a przy $p$ lokalny dysk $W^s(p)$.
Wyjaśnijmy wersję lematu Morse'a dostosowaną do takiego dysku.
Gradient jest do niego styczny. We współrzędnych
tubularnych opartych na normalnej ortogonalnej do dysku
różniczka funkcji w kierunku normalnym znika na dysku.
Taylor w zmiennej normalnej daje zatem
\[
 F(z,w)=F(z,0)+w^{\mathsf T}C(z,w)w.
\]
Przy dysku niestabilnym $C$ jest dodatnio określona,
a $F(z,0)$ ma niezdegenerowane maksimum.
Lemat Morse'a na dysku i zamiana $w\mapsto C(z,w)^{1/2}w$
dają standardowy model, zachowując dysk $w=0$.
Przy dysku stabilnym znaki są odwrotne.
W tych mapach zastępujemy pole standardowym polem
spadku, sklejonym ze starym za pomocą funkcji odcinającej.
Wypukłe połączenie dwóch pól z $XF<0$ nadal obniża $F$.
Oba pola są styczne do zachowywanego dysku, więc jego
ślad poza małym otoczeniem punktu krytycznego się nie zmienia.
W szczególności zachowujemy jedyne przecięcie $A_q\cap B_p$.

Modelem dla całej pary będzie
\begin{equation}\label{eq:narodziny-uchwytow-18}
 G_s(t,x,y)=\frac{t^3}{3}-st-|x|^2+|y|^2,
 \quad x\in\R^k,\quad y\in\R^{m-k-1}.
\end{equation}
Po dodatniej normalizacji afinicznej wartości funkcji
modele przy $p,q$ identyfikujemy z otoczeniami dwóch
punktów krytycznych $G_{s_0}$, gdzie $s_0>0$.
Stosując powyższe sklejenie, możemy tam od razu
wybrać pole odpowiadające $-\nabla G_{s_0}$.

\emph{Krok 3: uzgadniamy przejście przez regularny pas.}
Małe przekroje poniżej $q$ i powyżej $p$ łączy
dyfeomorfizm $h$ wyznaczony przepływem.
W modelu $G_{s_0}$ mamy analogiczny dyfeomorfizm $h_0$.
Po identyfikacji przekrojów ich różnicę opisuje lokalny
dyfeomorfizm $j$ przestrzeni $\R^k\times\R^{m-k-1}$,
ustalający $0$. Warunek transwersalności oznacza,
że $j(\R^k\times\{0\})$ jest poprzeczny do
$\{0\}\times\R^{m-k-1}$ w zerze.
Na małym dysku rodzina $j_\tau(z)=\tau^{-1}j(\tau z)$,
uzupełniona przez $j_0=Dj_0$, jest gładką izotopią
do części liniowej. Gładkość przy $\tau=0$ wynika ze wzoru
$j_\tau(z)=\int_0^1 Dj(t\tau z)z\,\dd t$.

Także dopasowanie części liniowej wymaga uzasadnienia.
Niech $Dj_0=\left(\begin{smallmatrix}A&B\\C&D\end{smallmatrix}\right)$.
Transwersalność daje odwracalność $A$, a odwracalność
$Dj_0$ daje odwracalność $S=D-CA^{-1}B$.
Mamy rozkład
\[
 Dj_0=
 \begin{pmatrix}I&0\\CA^{-1}&I\end{pmatrix}
 \begin{pmatrix}A&0\\0&S\end{pmatrix}
 \begin{pmatrix}I&A^{-1}B\\0&I\end{pmatrix}.
\]
Ścinania zewnętrzne kurczymy liniowo do identyczności.
Odbicia wybranych osi w mapach modelowych pozwalają
uzgodnić znaki wyznaczników $A,S$; nie zmieniają funkcji $G_s$.
Następnie łączymy oba bloki z identycznością w
$\operatorname{GL}^+$ (rozkład polarny i spójność grupy
obrotów). Cała droga zachowuje odwracalność bloku $A$,
a więc transwersalność. Bloki wymiaru zero pomijamy.

Izotopię lokalizujemy, mnożąc jej pole generujące
przez funkcję odcinającą. Kontrolujemy przy tym
\emph{całe} przecięcie, nie tylko różniczkę w zerze.
Na dostatecznie małym dysku, jednostajnie w parametrze,
\[
 \operatorname{dist}(j_\tau(u,0),\{0\}\times\R^{m-k-1})
       \geq c|u|,\qquad |V_\tau(z)|\leq C|z|.
\]
Pierwsza nierówność wynika z jednostajnej odwracalności
bloku poprzecznego, druga z $V_\tau(0)=0$.
Dzielimy izotopię na skończenie wiele dostatecznie
krótkich kroków. W każdym kroku odcięte pole zgadza się
z właściwym polem na mniejszym dysku i znika poza większym.
Czas kroku wybieramy tak, aby przesunięcie w pierścieniu
było mniejsze niż połowa dolnego oszacowania odległości.
Wtedy poza zerem nie powstaje żadne nowe przecięcie.
Kolejne mniejsze dyski wybieramy tak, aby ich obrazy
pozostawały w obszarze zgodności poprzedniego kroku.
Po skończonej liczbie kroków nowe odwzorowanie przejścia
zgadza się z $h_0$ blisko zera, a poza większym dyskiem
ze starym odwzorowaniem.

Realizujemy tę izotopię przez zmianę pola w regularnym paśmie.
We współrzędnej czasu $u=-F$ stare znormalizowane pole
jest $\partial_u$. Dodanie do niego pola stycznego
do przekrojów, generującego wybraną izotopię,
nie zmienia równości $XF=-1$.
Parametr izotopii wybieramy stały przy końcach pasa.
Nowe pole skleja się więc ze starym i nadal jest polem
spadku; jego odwzorowanie przejścia to wymagane $h_0$.
Przedłużając współrzędne przez przepływ, otrzymujemy
na otoczeniu całej trajektorii model $-\nabla G_{s_0}$
z dokładnością do dodatniej zmiany szybkości.
To wymagało zmiany pola, nie wynikało z samego
lematu Morse'a dla pierwotnego gradientu.

Dla ustalonego $s=s_0>0$ model ma dwa punkty:
$t=-\sqrt{s_0}$ o indeksie $k+1$
i $t=\sqrt{s_0}$ o indeksie $k$.
Wzdłuż osi $x=y=0$ ujemny gradient
spełnia $\dot t=s_0-t^2$ i biegnie
od pierwszego punktu do drugiego.
Każda inna trajektoria ma niezerową
składową $x$ albo $y$ i nie może mieć
obu tych końców; jest to jedyna
trajektoria łącząca.

Potrzebna jest jeszcze kontrola powrotów poza ten model.
Dla dostatecznie małego otoczenia $U$ domkniętej trajektorii
$\gamma\cup\{p,q\}$ istnieje mniejsze $V$, którego domknięcie
leży w $U$, takie że stara trajektoria wychodząca z $V$
i opuszczająca $U$ nie wraca do $V$.
Gdyby nie, istniałby ciąg odcinków wychodzących
z punktów dążących do tego łuku, przechodzących poza $U$
i ponownie zbliżających się do łuku.
Po wybraniu punktu granicznego poza $U$ otrzymalibyśmy
inną trajektorię od $q$ do $p$.
Dokładniej: gdyby jej koniec lub początek leżał
w brzegu pasa, gładka zależność i skończony czas
trafienia zmuszałyby do tego także pobliskie odcinki.
Nie mogłyby one mieć obu końców blisko wewnętrznego łuku.
Jedynymi możliwymi końcami krytycznymi są $q,p$,
a ścisły spadek wyklucza powrót do tego samego punktu.
Otrzymana druga trajektoria przeczy założeniu.
Argument stosujemy do pola \emph{po} uzgodnieniu modeli;
właśnie dlatego zachowaliśmy jedno przecięcie w całej izotopii.

\emph{Krok 4: usuwamy zera pola.}
Po pomnożeniu pola spadku przez
dodatnią funkcję jego trajektorie
w modelowej rurze są trajektoriami
\[
 X=(s_0-t^2,\,2x,\,-2y).
\]
Wybierzmy funkcję $v_1(t)<0$,
która zgadza się z $s_0-t^2$
blisko obu końców rury
($|t|>\sqrt{s_0}$), a różnicę podpieramy
w rzucie mniejszego otoczenia $V$. Niech
$\rho(r)=1$ dla bardzo małych
$r\geq0$ i $\rho(r)=0$ przy
brzegu małego przekroju, tak aby nośnik całej
zmiany był zawarty w $V$. Kładziemy
\[
 v(t,r)=\rho(r)v_1(t)
       +(1-\rho(r))(s_0-t^2),\qquad
 \widetilde X=
 \bigl(v(t,|x|^2+|y|^2),\,2x,\,-2y\bigr).
\]
Pole $\widetilde X$ jest równe $X$
przy brzegu rury i nigdzie nie
znika: dla $(x,y)\ne(0,0)$ jedna
ze składowych poprzecznych jest
niezerowa, a na osi $v=v_1<0$.
Przedłużamy je starym polem na cały pas.
Pole zmieniliśmy tylko w $V$, a wzory modelowe
obowiązują w większym $U$. Trajektoria nie może
pozostać w $U$ przez nieskończony czas
w przyszłości: wtedy z równania
$\dot x=2x$ wynika $x=0$,
z $\dot y=-2y$ wynika $y\to0$,
a na małym przekroju $v<0$
każe $t$ opuścić rurę w skończonym
czasie. W przeszłości argument
jest symetryczny, z rolami $x,y$
zamienionymi. Gdy opuści $U$, biegnie starym polem
aż do ewentualnego następnego trafienia w $V$.
Taki powrót jest wykluczony przez powyższy argument:
od ostatniego wyjścia z $V$ do opuszczenia $U$
również biegła starym polem. Poza $V$ nie może
też asymptotycznie zbiegać do $p$ lub $q$.
Każda trajektoria nowego pola biegnie więc od górnego
brzegu pasa do dolnego.

\emph{Krok 5: budujemy funkcję.}
Przepływ pola $-\widetilde X$
traktowanego jako skierowane ku
górze daje dyfeomorfizm pasa
z $N_{\rm dolny}\times[0,1]$.
Czas trafienia w górny brzeg jest skończony
i gładki dzięki poprzeczności. Dzieląc czas przez
ten czas trafienia na każdej trajektorii, otrzymujemy
parametr przedziału $[0,1]$ i gładki dyfeomorfizm.
Oznaczmy współrzędną przedziału przez $u$; pole $-\widetilde X$
jest dodatnią wielokrotnością
$\partial_u$. Na dwóch małych
kołnierzach stara funkcja $F$
ściśle rośnie wraz z $u$.
Na czas tej konstrukcji normalizujemy
afinicznie jej wartości na brzegach
do $0$ i $1$.
Wybieramy $\lambda(u)$ równą $1$
blisko $0$ i $1$, a równą $0$
na środku. Po wystarczającym
zmniejszeniu kołnierzy mamy
$\int_0^1\lambda\,\partial_uF\,du<1$.
Dla każdego $z\in N_{\rm dolny}$ niech
\[
 h(z)=
 \frac{1-\int_0^1\lambda(u)
                 \partial_uF(z,u)\,du}
      {\int_0^1(1-\lambda(u))\,du}>0,
 \quad
 \mu(z,u)=\lambda(u)\partial_uF(z,u)
            +(1-\lambda(u))h(z).
\]
Funkcja $\mu$ jest dodatnia i ma
całkę $1$ na każdej trajektorii.
Zatem $F'(z,u)=\int_0^u\mu(z,v)\,dv$
nie ma punktów krytycznych i przyjmuje
wartości $0,1$ na odpowiednich brzegach.
Ponieważ $\mu=\partial_uF$ blisko
obu końców, $F'=F$ w kołnierzach.
Na koniec cofamy normalizację afiniczną.
Uzyskaliśmy iloczynowy pas bez
zmiany na brzegu. Rodzina $G_s$
pokazuje na wykresie ten sam
mechanizm lokalny: dla $s<0$
jej pochodna $t^2-s$ jest
wszędzie dodatnia.
\end{proof}

\begin{figure}[htbp]
\centering
\begin{tikzpicture}[>=Latex,line cap=round,font=\small,scale=1.04]
 \draw[->,gray!70] (-2.0,0)--(2.0,0) node[right] {$t$};
 \draw[->,gray!70] (0,-1.32)--(0,1.32) node[above] {$G_s$};
 \draw[manifoldblue!80!black,very thick,
       domain=-1.64:1.64,samples=65]
      plot (\x,{\x*\x*\x/3-\x});
 \fill[manifoldred] (-1,{2/3}) circle (2.5pt)
       node[above left] {$q$};
 \fill[manifoldblue!85!black] (1,{-2/3}) circle (2.5pt)
       node[below right] {$p$};
 \node at (0,-1.72) {$s=1$: para punktów};
 \begin{scope}[shift={(5.3,0)}]
  \draw[->,gray!70] (-2.0,0)--(2.0,0) node[right] {$t$};
  \draw[->,gray!70] (0,-1.32)--(0,1.32) node[above] {$G_s$};
  \draw[manifoldgreen!80!black,very thick,
        domain=-1.20:1.20,samples=60]
       plot (\x,{\x*\x*\x/3+\x/2});
  \node at (0,-1.72) {$s=-1/2$: brak punktów};
 \end{scope}
\end{tikzpicture}
\caption{Przekrój modelu~\eqref{eq:narodziny-uchwytow-18}
przez oś $x=y=0$. Dla $s>0$ są dwa
punkty krytyczne o kolejnych indeksach
i jedna trajektoria od $q$ do $p$.
Po przejściu przez $s=0$ oba znikają.
Pozostałe kierunki $x,y$ ustalają indeksy.}
\label{fig:znoszenie-uchwytow-18}
\end{figure}

\begin{stwierdzenie}[Warunkowe przejście od przecięcia algebraicznego]
\label{prop:whitney-do-zniesienia-18}
Niech dla pary sąsiednich uchwytów
podpisana suma punktów $A_q\cap B_p$
wynosi $\pm1$. Jeśli punktów jest więcej niż jeden,
załóżmy, że można sparować wszystkie poza jednym
w pary o przeciwnych znakach i wybrać dla nich
rozłączne \emph{osadzone, prawidłowo obramowane
dyski Whitneya}. Brzeg każdego dysku składa się
z łuku na $A_q$ i łuku na $B_p$, a jego wnętrze
omija wszystkie rozważane sfery. Wnętrza obu łuków
omijają pozostałe punkty $A_q\cap B_p$.
Wtedy po izotopii sfer
pozostaje jeden punkt przecięcia
i stosuje się
Twierdzenie~\ref{thm:zniesienie-18}.
\end{stwierdzenie}
\begin{proof}
Jeśli jest już jeden punkt, od razu stosujemy
twierdzenie o zniesieniu. W pozostałym przypadku
założenie o łukach wymaga dodatnich wymiarów obu sfer;
w szczególności $m\geq3$ i $1\leq k\leq m-2$,
gdzie indeksy uchwytów to $k,k+1$.
Wybierzmy jeden z założonych dysków.
„Prawidłowo obramowany” oznacza tutaj, że obramowanie
wzdłuż jego brzegu, wyznaczone przez kierunki
styczne i normalne obu sfer, przedłuża się na dysk.
Nie wystarcza sama trywialność jego wiązki normalnej:
ta dla dysku jest automatyczna, a warunek dotyczy
\emph{zadanego} obramowania brzegowego.
Rozpiszmy ruch w tym otoczeniu. Niech $r=k$ i $s=m-k-1$
będą wymiarami $A_q$ i $B_p$. Zgodne obramowanie i otoczenie
tubularne dają produktowy model
$U\times D^{r-1}\times D^{s-1}$, gdzie $U$ jest małym
otoczeniem płaskiego dysku o brzegu złożonym z dwóch łuków.
W tym modelu fragmenty sfer mają postać
\[
 A_q=\alpha\times D^{r-1}\times\{0\},\qquad
 B_p=\beta\times\{0\}\times D^{s-1}.
\]
Łuki $\alpha,\beta$ przedłużamy nieco poza ich dwa przecięcia.
Obramowanie uzgadnia kierunki styczne do sfer wzdłuż tych łuków
z kierunkami normalnymi do dysku; jego przedłużenie na dysk
pozwala pogrubić cały model, a nie tylko jego brzeg.
Jest to lokalny model z lematu~6.7 Milnora, tu przy założonym
istnieniu osadzonego, czystego i prawidłowo obramowanego dysku.

W płaszczyźnie wybieramy izotopię $G_t$ przesuwającą $\alpha$
przez dysk na drugą stronę $\beta$, tak aby $G_1(\alpha)$
nie przecinała $\beta$. Jej nośnik mieści się we wnętrzu $U$.
Niech $\rho(u,v)$ będzie funkcją odcinającą równą $1$ blisko
$(0,0)$ i $0$ przy brzegu czynników normalnych. Wówczas
\[
 H_t(z,u,v)=\bigl(G_{t\rho(u,v)}(z),u,v\bigr)
\]
jest izotopią otoczenia, równą identyczności blisko jego brzegu:
dla ustalonych $u,v$ odwrotność daje $G_{t\rho(u,v)}^{-1}$.
Przecięcie $H_t(A_q)$ z nieruchomą $B_p$ wymaga $u=v=0$,
więc odpowiada dokładnie przecięciom $G_t(\alpha)$ z $\beta$.
Dla $t=1$ oba punkty znikają. Czystość dysku, warunek o wnętrzach
łuków i rozłączność dysków pozwalają wybrać otoczenie omijające
pozostałe punkty przecięcia i pozostałe dyski. Poza nim kładziemy
identyczność. Normalne pogrubienie przesuwanej sfery przenosimy
przez tę samą izotopię, zachowując dane przyczepienia uchwytu.
Powtarzamy ten ruch dla rozłącznych
dysków. Po każdym kroku liczba
punktów zmniejsza się o dwa, a
suma podpisana nie zmienia się.
Ponieważ na końcu wynosi $\pm1$,
pozostaje jeden punkt; poprzednie
twierdzenie usuwa parę uchwytów.
\end{proof}

\begin{uwaga}[Dlaczego potrzebny będzie wymiar]
Samo równanie $I(A_q,B_p)=\pm1$
nie zapewnia istnienia wymienionych
dysków. Trzeba znaleźć odpowiednie
łuki, wypełnić ich pętle, odsunąć
wnętrza dysków od sfer i sprawdzić
obramowanie. W wysokim wymiarze
pomagają prosta spójność i pozycja
ogólna; przy indeksach brzegowych
oraz w małych wymiarach występują
dodatkowe trudności. Szczegółowe
warunki triku Whitneya są zasadniczą
częścią dowodu twierdzenia Smale'a,
a nie konsekwencją samego rachunku
macierzy.
\end{uwaga}

\section{Granica tego rozdziału i dalszy krok}
\label{sec:uchwyty-smale-most-18}

Rozkład uchwytowy daje skończony
kompleks względny. Jeśli w kobordyzmie
obie inkluzje końców są równoważnościami
homotopijnymi, mówimy o
\emph{$h$-kobordyzmie}. Wtedy
$H_*(W,N_-)=0$, więc kompleks
uchwytów jest acykliczny.
Rachunek macierzy sugeruje pary
uchwytów do zniesienia, lecz
aby naprawdę uzyskać iloczyn
$N_-\times[0,1]$, trzeba
zrealizować operacje algebraiczne
przez przesunięcia i usunąć
zbędne przecięcia trikiem Whitneya.
W przypadku prostej spójności i
wymiaru $m\geq6$ prowadzi to do
twierdzenia Smale'a o $h$-kobordyzmie.
Bez prostej spójności pojawia się
dodatkowa przeszkoda, torsja
Whiteheada; jej analiza należy
do twierdzenia o $s$-kobordyzmie.
Nie jest to inwariant samego powyższego kompleksu nad $\Z$.
Dla spójnego $W$ i grupy $\pi=\pi_1(W)$ trzeba podnieść uchwyty
do nakrycia uniwersalnego i pracować nad pierścieniem
grupowym $\Z[\pi]$. Współczynniki zapamiętują wtedy
także klasy dróg, a nie tylko sumy znaków przecięć.
Acykliczność zwykłego kompleksu całkowitego
nie wystarcza do rozpoznania tej przeszkody.

Drugi kierunek to teoria chirurgii:
Twierdzenie~\ref{thm:slad-chirurgii-18}
tłumaczy każdy uchwyt na operację
na poziomicy, a obramowanie mówi,
jak wykonać sklejanie. Naturalnym dalszym krokiem byłoby zbadanie, jakie
inwarianty homologiczne i formy
przecięcia można zmienić przez
dobór takich operacji.
Nie zakładamy tu, że wszystkie
przeszkody można usunąć;
zwłaszcza przypadki małych
wymiarów wymagają osobnej analizy.

\par\smallskip
{\footnotesize\noindent\emph{Dalsza lektura:}
\href{https://www.cambridge.org/core/books/differential-topology/theory-of-handle-decompositions/BA54EF1BD9B755DFA2B363176B0CA2F7}{C.\,T.\,C.\ Wall, \emph{Differential Topology}, rozdział o rozkładach uchwytowych},
\href{https://webhomes.maths.ed.ac.uk/~v1ranick/books/surgery.pdf}{A.\ Ranicki, \emph{Algebraic and Geometric Surgery}, ślad chirurgii i kompleks uchwytowy},
\href{https://www.math.columbia.edu/~jmorgan/Lecture_IIIA_hcobordism_Contd.pdf}{J.\ Morgan, wykład o znoszeniu uchwytów i triku Whitneya},
\href{https://webhomes.maths.ed.ac.uk/~v1ranick/surgery/hcobord.pdf}{J.\ Milnor, \emph{Lectures on the h-Cobordism Theorem}, \S6, zwłaszcza lemat~6.7 i ruch w jego produktowym modelu}.\par}
